
This work presents empirical measurements of k-means clustering performance on epoch-5 penultimate-layer embeddings of a pretrained ResNet-50, applied to Waterbirds and CelebA. The same procedure yields AMI=0.762 and purity=0.892 on Waterbirds, and AMI=0.258 and purity=0.456 on CelebA. Both datasets exhibit spurious correlations under ERM; the detection procedure succeeds on one and fails on the other under identical experimental conditions.

The analysis attributes this discrepancy to the alignment between each dataset's spurious attribute and the ImageNet pretraining prior. Waterbirds' spurious attribute (scene-level background habitat) is well-represented in ImageNet-pretrained features, producing cluster-separable embedding structure by epoch 5. CelebA's spurious attribute (hair color texture) is not represented in the pretraining prior at the required level of separability, and five epochs of ERM fine-tuning do not establish this structure.

The characterization of this conditionality — that early-epoch annotation-free clustering detection depends on pretraining alignment rather than on spurious correlation strength in the downstream task — represents the main finding of this work (C1). The empirical gap is quantified (C2), a candidate diagnostic is proposed with the caveat that empirical calibration is needed (C3), and a configuration-sensitivity reporting gap in the published literature is identified (C4).

The downstream GSB intervention remains unvalidated because the CelebA detection prerequisite was not satisfied. Three directions that follow from these results are: adaptive probe epoch selection based on representational stability (addressing L3), k-adaptive clustering for multi-group spurious structures (addressing L2), and Waterbirds-restricted evaluation of the gradient SNR mechanism (addressing L4) where detection is confirmed to be reliable.

